% ============================================================
% Linear Control Systems HW3 Report Template (EMPTY)
% Obeys: <= 8 pages main report; Appendix not counted
% ============================================================

\documentclass[11pt]{article}

\usepackage[margin=1in]{geometry}
\usepackage{amsmath,amssymb}
\usepackage{booktabs}
\usepackage{graphicx}
\usepackage{hyperref}
\usepackage{enumitem}
\usepackage{caption}
\usepackage{subcaption}

% --------- Metadata (edit these) ----------
\newcommand{\CourseName}{Project Proposal}
\newcommand{\HWName}{Side-Jet Interactions over a Hypersonic Cone Flow}
\newcommand{\StudentName}{Mehmet Furkan Tunc, Elif Avcılar, Bayram Celik}

\newcommand{\DateSubmitted}{\today}

% --------- Convenience ----------
\newcommand{\tol}{\texttt{tol}}
\newcommand{\MATLAB}{\textsc{Matlab}}
\newcommand{\RKFour}{RK4}

\begin{document}

% ============================================================
% Title Page (Page 1)
% ============================================================
\begin{titlepage}
\centering

\vspace*{1.0in}

{\LARGE \CourseName \par}
\vspace{0.15in}
{\Large \HWName \par}

\vspace{0.75in}

{\large \textbf{\StudentName} \par}


\vspace{0.5in}

\vfill

{\large \DateSubmitted \par}

\end{titlepage}

% ============================================================
% Methods (<= 1 page)
% ============================================================
\section{Motivation}
Side-jet reaction control is widely used for maneuvering and reentry at high Mach numbers, yet the resulting jet–crossflow interaction can introduce separation/reattachment, multiple recirculation zones, and strong shock/shock and shock/shear-layer coupling that substantially alters surface pressure, aeroheating, and integrated forces (i.e., aerodynamic interference). 

In this context, Karpuzcu and Levin (2022, 2023) conducted detailed hypersonic side-jet studies—examining canonical geometries including cones and reporting complex flow structures and unsteadiness using kinetic simulations (and, in 2022, complementary wind-tunnel schlieren evidence). 

However, their baseline configurations emphasize cold, non-reacting conditions, and therefore do not address thermochemical non-equilibrium (e.g., vibrational excitation/dissociation) that becomes relevant for higher-enthalpy hypersonic environments. 

Moreover, these works do not provide a direct, quantitative comparison against hypersonic CFD in terms of both predictive accuracy and computational cost—precisely the trade space needed for practical parametric design studies. 

Finally, while side-jet interactions are frequently simulated with continuum CFD, much of that numerical literature targets supersonic crossflow regimes, leaving a narrower CFD evidence base for hypersonic cone side-jet conditions.

In this study, our primary objectives are:
\begin{itemize}
    \item to study non-equilibrium effects on the flowfield characteristics using DSMC method in different configurations and,
    \item testing hypersonic CFD solver capabilities against experimental and numerical results in terms of accuracy and cost using the same configuration as in base line study.
\end{itemize}  

\section{Methodology and Configuration}
In this study, following solvers will be employed:
\begin{itemize}
    \item \texttt{dsmcFoam+} is an OpenFOAM-based Direct Simulation Monte Carlo (DSMC) solver used to compute rarefied/kinetic jet--crossflow interactions with time-accurate sampling. 
    \item \texttt{hy2Foam} is an OpenFOAM finite-volume hypersonic CFD solver (two-temperature framework) used for continuum-based comparisons of accuracy and computational cost.
\end{itemize}
Karpuzcu and Levin (2023) report that a single DSMC simulation advanced to a final physical time of $t_f=5~\mathrm{ms}$ required approximately $4.6\times 10^{5}$ CPU-hours on Intel Xeon Phi 7250 ("Knights Landing") processors. In the present work, the planned thermochemical non-equilibrium DSMC campaign (additional internal-energy/chemistry modeling, tighter resolution requirements, and multiple operating points) is expected to demand significantly higher computational resources plus the hypersonic CFD comparison study which also requires resources (estimating $2.0\times 10^{6}$ CPU-hours for the complete non-equilibrium DSMC study set and hypersonic CFD comparison study).

Base-line study inputs are given in Table \ref{tab:baselineInputs}. Relative to base-line study, differences of planned case studies are given in Table \ref{tab:differences}.
\begin{table}[t]
\centering
\small
\setlength{\tabcolsep}{6pt}
\renewcommand{\arraystretch}{1.15}
\begin{tabular}{l l}
\hline
\multicolumn{2}{c}{\textbf{Geometry and placement (normalized by cone length $L$)}}\\
\hline
Cone diameter & $D = 0.10~\mathrm{m}$ \\
Cone length & $L = 0.40~\mathrm{m}$ \\
Jet diameter & $d_j = 1.0\times 10^{-3}~\mathrm{m}$ \\
Jet center location & $x/L = 0.50,\;\; y/L = 0.0625$ \\
Freestream unit Reynolds number & $Re_\infty = 1.0\times 10^{5}~\mathrm{m^{-1}}$ \\
\hline
\multicolumn{2}{c}{\textbf{DSMC time-accurate settings}}\\
\hline
Time step & $\Delta t = 5.0\times 10^{-9}~\mathrm{s}$ \\
Flow establishment time & $t_{\mathrm{est}} = 2.5\times 10^{-3}~\mathrm{s}$ \\
Sampling window & $t_{\mathrm{samp}} = 5.0\times 10^{-4}~\mathrm{s}$ \\
Final physical time & $t_f = 5.0\times 10^{-3}~\mathrm{s}$ \\
\hline
\multicolumn{2}{c}{\textbf{Molecular/physical modeling (baseline cold case)}}\\
\hline
Molecular model & Variable Hard Sphere (VHS) \\
$\mathrm{N_2}$ molecular diameter & $d_{\mathrm{N_2}} = 4.467\times 10^{-10}~\mathrm{m}$ \\
Viscosity index & $\omega = 0.26$ \\
Wall boundary condition & Isothermal, $T_w = 300~\mathrm{K}$ \\
Chemistry & None (no reactions) \\
Dimensionality & Axisymmetric \\
Working fluids & Freestream: $\mathrm{N_2}$,\;\; Jet: $\mathrm{N_2}$ \\
\hline
\end{tabular}
\caption{Baseline study inputs}
\label{tab:baselineInputs}
\end{table}

\begin{table}[t]
\centering
\small
\setlength{\tabcolsep}{6pt}
\renewcommand{\arraystretch}{1.15}
\begin{tabular}{l l l l l}
\hline
\multicolumn{5}{c}{\textbf{DSMC Studies}}\\
\hline
Case Number & Gas & Thermal Eq. & Chemical Eq. & Method\\
\hline
Baseline Study & $N_2$ & Yes & Yes & DSMC \\
Case 1 & $N_2$ & No & Yes & DSMC\\
Case 2 & $N_2$ & No & No & DSMC\\
Case 3 & 0.8 $N_2$ - 0.2 $O_2$ & Yes & Yes & DSMC\\
Case 4 & 0.8 $N_2$ - 0.2 $O_2$ & No & No & DSMC\\
\hline
\multicolumn{5}{c}{\textbf{CFD Studies}}\\
\hline
Case Number & Gas & Thermal Eq. & Chemical Eq. & Method\\
\hline
Baseline Study & $N_2$ & Yes & Yes & DSMC \\
Case 1 & $N_2$ & Yes & Yes & CFD\\
Case 2 & $N_2$ & No & No & CFD\\
Case 3 & 0.8 $N_2$ - 0.2 $O_2$ & Yes & Yes & CFD\\
\hline
\end{tabular}
\caption{Differences in literature and proposed studies}
\label{tab:differences}
\end{table}

\end{document}

